\subsection{Propiedades funtoriales de los invariantes normales de una inmersión}
\label{IV.16.2}

\begin{env}[16.2.1]
\label{IV.16.2.1}
Sean $f=(\psi,\theta):Y\to X$,
$f'=(\psi',\theta'):Y'\to X'$ dos morfismos de espacios anillados tales que los
homomorfismos $\theta^\#$ y $\theta'^\#$ sean sobreyectivos; consideremos un
diagrama conmutativo de morfismos de espacios anillados
\[
  \xymatrix{
    Y \ar[r]^f & X \\
    Y' \ar[r]_{f'} \ar[u]^u & X' \ar[u]_v
  }
  \tag{16.2.1.1}
  \label{IV.16.2.1.1}
\]
Pongamos $u=(\rho,\lambda)$, $v=(\sigma,\mu)$. Se tiene
$\rho^*(\psi^*(\mathcal{O}_X))=
\psi'^*(\sigma^*(\mathcal{O}_X))$ y, por consiguiente, un diagrama conmutativo
de homomorfismos de haces de anillos sobre $Y'$
\[
  \xymatrix{
    \rho^*(\psi^*(\mathcal{O}_X))
      =\psi'^*(\sigma^*(\mathcal{O}_X))
      \ar[r]^-{\psi'^*(\mu^\#)}
      \ar[d]_{\rho^*(\theta^\#)}
      & \psi'^*(\mathcal{O}_{X'}) \ar[d]^{\theta'^\#} \\
    \rho^*(\mathcal{O}_Y) \ar[r]_-{\lambda^\#}
      & \mathcal{O}_{Y'}
  }
\]
de donde se concluye, si $\mathscr{I}$ y $\mathscr{I}'$ son los núcleos de
$\theta^\#$ y $\theta'^\#$, que se tiene
$\psi'^*(\mu^\#)(\rho^*(\mathscr{I}))\subset\mathscr{I}'$, teniendo en cuenta la
exactitud del funtor $\rho^*$. Se deduce inmediatamente que, para todo entero
$n$,
$\psi'^*(\mu^\#)(\rho^*(\mathscr{I}^n))\subset\mathscr{I}'^n$, lo que muestra
que $\psi'^*(\mu^\#)$ define, al pasar a los cocientes, un homomorfismo de
haces de anillos
\[
  \nu_n:
  \rho^*(\psi^*(\mathcal{O}_X)/\mathscr{I}^{n+1})
  \longrightarrow
  \psi'^*(\mathcal{O}_{X'})/\mathscr{I}'^{n+1}
  \tag{16.2.1.2}
  \label{IV.16.2.1.2}
\]
y, por consiguiente, un morfismo de espacios anillados
$w_n=(\rho,\nu_n):Y'^{(n)}\to Y^{(n)}$ (que, para $n=0$, no es sino
$u$). De esta definición se sigue inmediatamente que los diagramas
\[
  \xymatrix@R=1pc{
    Y^{(n)} \ar[r]^-{h_{mn}} & Y^{(m)} \ar[r]^-{h_m} & X \\
    & & & (n\leq m) \\
    Y'^{(n)} \ar[r]_-{h'_{mn}} \ar[uu]^-{w_n}
      & Y'^{(m)} \ar[r]_-{h'_m} \ar[uu]^-{w_m}
      & X' \ar[uu]_-v
  }
\]
(donde las flechas horizontales son los morfismos canónicos (16.1.2)) son
conmutativos.

Al pasar a los cocientes a partir de los homomorfismos (16.2.1.2), y teniendo en
\oldpage[IV-4]{10}
cuenta la exactitud del funtor $\rho^*$, se obtiene un dihomomorfismo
de álgebras graduadas (relativo al homomorfismo
$\lambda^\#:\rho^*(\mathcal{O}_Y)\to\mathcal{O}_{Y'}$)
\[
  \operatorname{gr}(u):
  \rho^*(\mathscr{G}r_\bullet(f))
  \longrightarrow\mathscr{G}r_\bullet(f')
  \tag{16.2.1.3}
  \label{IV.16.2.1.3}
\]
(o, si se prefiere, un $\rho$-morfismo (0\textsubscript{I}, 3.5.1)
$\mathscr{G}r_\bullet(f)\to\mathscr{G}r_\bullet(f')$), y en particular un
dihomomorfismo de los haces conormales
\[
  \operatorname{gr}_1(u):
  \rho^*(\mathscr{G}r_1(f))\longrightarrow\mathscr{G}r_1(f').
\]
Es además inmediato que estos homomorfismos dan lugar a un diagrama
conmutativo
\[
  \xymatrix{
    \rho^*(\mathscr{S}_{\mathcal{O}_Y}^{\bullet}
      (\mathscr{G}r_1(f)))
      \ar[r] \ar[d]_-{\mathscr{S}(\operatorname{gr}_1(u))}
      & \rho^*(\mathscr{G}r_\bullet(f))
        \ar[d]^-{\operatorname{gr}(u)} \\
    \mathscr{S}_{\mathcal{O}_{Y'}}^{\bullet}
      (\mathscr{G}r_1(f'))
      \ar[r]
      & \mathscr{G}r_\bullet(f')
  }
  \tag{16.2.1.4}
  \label{IV.16.2.1.4}
\]
donde las flechas horizontales son los homomorfismos canónicos (16.1.2.2).

Finalmente, si se tiene un diagrama conmutativo de morfismos de espacios anillados
\[
  \xymatrix{
    Y \ar[r]^f & X \\
    Y' \ar[r]_{f'} \ar[u]^u & X' \ar[u]_v \\
    Y'' \ar[r]_{f''} \ar[u]^{u'} & X'' \ar[u]_{v'}
  }
\]
donde $f''=(\psi'',\theta'')$ es tal que $\theta''^\#$ es sobreyectivo, y si
$w'_n$ y $w''_n$ se definen a partir de $u'$, $v'$, por una parte, y de
$u''=u\circ u'$, $v''=v\circ v'$, por otra, entonces se tiene
$w''_n=w_n\circ w'_n$, como se sigue inmediatamente de las definiciones y de
(0\textsubscript{I}, 3.5.5); se tiene asimismo
$\operatorname{gr}(u'')=
\operatorname{gr}(u')\circ\rho'^*(\operatorname{gr}(u))$ si
$u'=(\rho',\lambda')$. Se puede, por tanto, decir que los $Y^{(n)}$ y
$\mathscr{G}r_\bullet(f)$ dependen funtorialmente de $f$.
\end{env}

\begin{proposition}[16.2.2]
\label{IV.16.2.2}
Con las notaciones e hipótesis de (16.2.1), supongamos además que $f$, $f'$,
$u$ y $v$ son morfismos de preesquemas. Entonces:
\begin{enumerate}
\item[(i)] Los morfismos $w_n:Y'^{(n)}\to Y^{(n)}$ son morfismos de
preesquemas.
\item[(ii)] Si $Y'=Y\times_X X'$, siendo $u$ y $f'$ las proyecciones
canónicas, y si $f$ es una inmersión o si $v$ es plano, se tiene
$Y'^{(n)}=Y^{(n)}\times_X X'$.
\item[(iii)] Si $Y'=Y\times_X X'$ y si $v$ es plano (resp. si $f$ es una
inmersión), el homomorfismo
\[
  \operatorname{Gr}(u)=\operatorname{gr}(u)\otimes 1:
  \mathscr{G}r_\bullet(f)\otimes_{\mathcal{O}_Y}\mathcal{O}_{Y'}
  \longrightarrow\mathscr{G}r_\bullet(f')
\]
es biyectivo (resp. sobreyectivo).
\end{enumerate}
\end{proposition}

\noindent(i) Las hipótesis implican inmediatamente que, para todo $y'\in Y'$,
$\rho_{y'}^*(\theta_{\psi'(y')}^\#)$ es un homomorfismo local (I, 1.6.2),
por tanto $w_n$ es un morfismo de preesquemas (I, 2.2.1).

\oldpage[IV-4]{11}
\noindent(ii) y (iii) Si $f$ es una inmersión, basta considerar el caso en que
$f$ es una inmersión cerrada, estando entonces $Y$ definido por el ideal
cuasicoherente $\mathscr{I}$ de $\mathcal{O}_X$, y $Y^{(n)}$ por el ideal
$\mathscr{I}^{n+1}$; las afirmaciones resultan entonces de (I, 4.4.5).

Supongamos, en segundo lugar, que $v$ es plano; basta considerar el caso en que
$X=\operatorname{Spec}(A)$, $Y=\operatorname{Spec}(B)$,
$X'=\operatorname{Spec}(A')$ son afines, siendo $A'$ un $A$-módulo plano;
entonces $Y'=\operatorname{Spec}(B')$ con $B'=B\otimes_A A'$; además, si
$\mathfrak{I}$ es el núcleo del homomorfismo $A\to B$, el núcleo
$\mathfrak{I}'$ de $A'\to B'$ se identifica con $\mathfrak{I}\otimes_A A'$ por
planitud, y
\[
  \mathfrak{I}'^n/\mathfrak{I}'^{n+1}
  =(\mathfrak{I}^n/\mathfrak{I}^{n+1})\otimes_A A'.
\]
Se deduce inmediatamente, teniendo en cuenta (0\textsubscript{I}, 4.3.3), que el
$\mathcal{O}_{Y'}$-módulo
$\mathscr{I}'^n/\mathscr{I}'^{n+1}$ es igual a
\begin{align*}
&\psi'^*\!\left(\sigma^*
  \left((\mathfrak{I}^n/\mathfrak{I}^{n+1})^{\sim}\right)\right)
  \otimes_{\sigma^*(\mathcal{O}_X)}\mathcal{O}_{X'} \\
={}&
\psi'^*\!\left(\sigma^*
  \left((\mathfrak{I}^n/\mathfrak{I}^{n+1})^{\sim}\right)\right)
  \otimes_{\psi'^*(\sigma^*(\mathcal{O}_X))}
  \psi'^*(\mathcal{O}_{X'}) \\
={}&
\rho^*(\mathscr{I}^n/\mathscr{I}^{n+1})
  \otimes_{\rho^*(\psi^*(\mathcal{O}_X))}
  \psi'^*(\mathcal{O}_{X'}),
\end{align*}
y como, en particular, para $n=0$, se tiene
\[
  \mathcal{O}_{Y'}=
  \rho^*(\mathcal{O}_Y)
  \otimes_{\rho^*(\psi^*(\mathcal{O}_X))}
  \psi'^*(\mathcal{O}_{X'}),
\]
se obtiene un isomorfismo canónico de
$\mathscr{I}'^n/\mathscr{I}'^{n+1}$ con
\[
  \rho^*(\mathscr{I}^n/\mathscr{I}^{n+1})
  \otimes_{\rho^*(\mathcal{O}_Y)}\mathcal{O}_{Y'}
  =
  (\mathscr{I}^n/\mathscr{I}^{n+1})
  \otimes_{\mathcal{O}_Y}\mathcal{O}_{Y'},
\]
lo que prueba (iii). Pongamos ahora
$C_n=\Gamma(Y,\mathcal{O}_{Y^{(n)}})$,
$C'_n=\Gamma(Y',\mathcal{O}_{Y'^{(n)}})$. Puesto que $Y^{(n)}$ y
$Y'^{(n)}$ son esquemas afines (16.1.5), el núcleo
$\mathfrak{K}_n$ (resp. $\mathfrak{K}'_n$) del homomorfismo
$C_n\to C_{n-1}$ (resp. $C'_n\to C'_{n-1}$) es
$\Gamma(Y,\mathscr{I}^n/\mathscr{I}^{n+1})$ (resp.
$\Gamma(Y',\mathscr{I}'^n/\mathscr{I}'^{n+1})$); por tanto, de lo que
precede se deduce que $\mathfrak{K}'_n=\mathfrak{K}_n\otimes_A A'$. Ahora bien, se tiene un
diagrama conmutativo
\[
  \xymatrix{
    0 \ar[r] &
    \mathfrak{K}_n\otimes_A A' \ar[r]\ar[d]^r &
    C_n\otimes_A A' \ar[r]\ar[d]^{s_n} &
    C_{n-1}\otimes_A A' \ar[r]\ar[d]^{s_{n-1}} & 0 \\
    0 \ar[r] &
    \mathfrak{K}'_n \ar[r] &
    C'_n \ar[r] &
    C'_{n-1} \ar[r] & 0
  }
\]
donde la flecha vertical de la izquierda es biyectiva y las dos filas son exactas
(siendo $A'$ un $A$-módulo plano). Se deduce por inducción que $s_n$ es
biyectivo para todo $n$: esto es cierto por hipótesis para $n=0$, y el caso de un
$n$ cualquiera se deduce aplicando el lema de los cinco. Esto prueba la segunda
afirmación de (ii).

\begin{corollary}[16.2.3]
\label{IV.16.2.3}
Sean $g:X\to Y$, $u:Y'\to Y$ dos morfismos de preesquemas,
$X'=X\times_Y Y'$, $g':X'\to Y'$ y $v:X'\to X$ las proyecciones canónicas.
Sea $f:Y\to X$ una sección sobre $Y$ de $X$ (por tanto, una inmersión),
$f'=f_{(Y')}:Y'\to X'$ la sección sobre $Y'$ de $X'$ deducida de $f$ por el
cambio de base $u$. Entonces:
\begin{enumerate}
\item[(i)] El morfismo
$w_n:Y_{f'}'^{(n)}\to Y_f^{(n)}$ correspondiente a $f$, $f'$, $u$, $v$
(16.2.1) y el morfismo canónico
$h'_n:Y_{f'}'^{(n)}\to X'$ identifican $Y_{f'}'^{(n)}$ con el producto fibrado
$Y_f^{(n)}\times_X X'$.
\item[(ii)] Si se dota a $\mathcal{O}_{Y_f^{(n)}}$ (resp.
$\mathcal{O}_{Y_{f'}'^{(n)}}$) de la estructura de $\mathcal{O}_Y$-álgebra
definida por $g$ (resp. de la estructura de $\mathcal{O}_{Y'}$-álgebra definida
por $g'$) (16.1.6), el homomorfismo de $\mathcal{O}_{Y'}$-álgebras
\[
  \rho^*(\mathcal{O}_{Y_f^{(n)}})
  \otimes_{\mathcal{O}_Y}\mathcal{O}_{Y'}
  \longrightarrow
  \mathcal{O}_{Y_{f'}'^{(n)}}
  \tag{16.2.3.1}
  \label{IV.16.2.3.1}
\]
\oldpage[IV-4]{12}
deducido del homomorfismo $\nu_n$ (16.2.1.2) es biyectivo. Además,
el homomorfismo de $\mathcal{O}_{Y'}$-módulos
\[
  \operatorname{Gr}_1(u):
  \mathscr{G}r_1(f)\otimes_{\mathcal{O}_Y}\mathcal{O}_{Y'}
  \longrightarrow\mathscr{G}r_1(f')
  \tag{16.2.3.2}
  \label{IV.16.2.3.2}
\]
es biyectivo.
\end{enumerate}
\end{corollary}

\noindent(i) Observemos primero que los morfismos $f':Y'\to X'$ y $u:Y'\to Y$
identifican $Y'$ con el producto fibrado $Y\times_X X'$ (para los morfismos estructurales
$f:Y\to X$ y $v:X'\to X$) (14.5.12.1). La conclusión de (i) resulta entonces de
(16.2.2, (ii)), pues el morfismo $g$ es una inmersión.

\noindent(ii) El diagrama conmutativo
\[
  \xymatrix{
    Y_f^{(n)} \ar[d]^{h_n}
      & Y_{f'}'^{(n)} \ar[d]^-{h'_n} \ar[l]^{w_n} \\
    X \ar[d]^g
      & X' \ar[d]^{g'} \ar[l]^v \\
    Y
      & Y' \ar[l]^u
  }
\]
identifica $Y_{f'}'^{(n)}$ con el producto fibrado $Y_f^{(n)}\times_X X'$ y $X'$ con el producto fibrado
$X\times_Y Y'$; por tanto, (I, 3.3.9) identifica (para los morfismos
$g'\circ h'_n$ y $w_n$) $Y_{f'}'^{(n)}$ con el producto fibrado
$Y_f^{(n)}\times_Y Y'$. Como $Y_f^{(n)}$ (resp. $Y_{f'}'^{(n)}$) es el
preesquema afín sobre $Y$ (resp. $Y'$) asociado al
$\mathcal{O}_Y$-álgebra $\mathcal{O}_{Y_f^{(n)}}$ (resp. al
$\mathcal{O}_{Y'}$-álgebra $\mathcal{O}_{Y_{f'}'^{(n)}}$), la biyectividad del
homomorfismo canónico (16.2.3.1) resulta de (II, 1.5.2).
Finalmente, el homomorfismo canónico (16.2.3.1) es compatible con las
aumentaciones
$\mathcal{O}_{Y_f^{(n)}}\to\mathcal{O}_Y$ y
$\mathcal{O}_{Y_{f'}'^{(n)}}\to\mathcal{O}_{Y'}$; como
$\mathcal{O}_{Y_f^{(n)}}$ es suma directa (en tanto que
$\mathcal{O}_Y$-módulo) de $\mathcal{O}_Y$ y del ideal de aumentación
$\mathscr{I}/\mathscr{I}^{n+1}$, se ve por tanto que el homomorfismo canónico
(16.2.3.1), restringido a
$(\mathscr{I}/\mathscr{I}^{n+1})\otimes_{\mathcal{O}_Y}\mathcal{O}_{Y'}$,
es una biyección de este último sobre
$\mathscr{I}'/\mathscr{I}'^{n+1}$. Para $n=1$, esto muestra que
$\operatorname{Gr}_1(u)$ es biyectivo.

Obsérvese que, bajo las hipótesis de (16.2.3), los homomorfismos
$\operatorname{Gr}_n(u)$ son sobreyectivos en virtud de lo que precede, pero no
son biyectivos en general para $n\geq 2$. Sin embargo:

\begin{corollary}[16.2.4]
\label{IV.16.2.4}
Bajo las hipótesis de (16.2.3), supongamos que $u:Y'\to Y$ es un morfismo
plano (resp. que los $\mathscr{G}r_n(f)$ son $\mathcal{O}_Y$-módulos
planos para $n\leq m$). Entonces el homomorfismo
\[
  \operatorname{Gr}_n(u):
  \mathscr{G}r_n(f)\otimes_{\mathcal{O}_Y}\mathcal{O}_{Y'}
  \longrightarrow\mathscr{G}r_n(f')
\]
es biyectivo para todo $n$ (resp. para $n\leq m$).
\end{corollary}

Si $u$ es plano, lo mismo ocurre con $v:X'\to X$, que se obtiene de él por
cambio de base, y ya se sabe en este caso que $\operatorname{Gr}(u)$ es
biyectivo (16.2.2, (iii)). Si los $\mathscr{G}r_n(f)$ son planos para
$n\leq m$, se ve primero por inducción sobre $n$ que lo mismo ocurre con los
$\mathscr{I}/\mathscr{I}^{n+1}$ para $n\leq m$, en virtud de las sucesiones exactas
\[
  0\longrightarrow\mathscr{I}^n/\mathscr{I}^{n+1}
  \longrightarrow\mathscr{I}/\mathscr{I}^{n+1}
  \longrightarrow\mathscr{I}/\mathscr{I}^n
  \longrightarrow 0
\]
\oldpage[IV-4]{13}
(0\textsubscript{I}, 6.1.2); además, se tienen entonces diagramas conmutativos
\[
  \xymatrix{
    0 \ar[r] &
    (\mathscr{I}^n/\mathscr{I}^{n+1})
      \otimes_{\mathcal{O}_Y}\mathcal{O}_{Y'}
      \ar[r]\ar[d] &
    (\mathscr{I}/\mathscr{I}^{n+1})
      \otimes_{\mathcal{O}_Y}\mathcal{O}_{Y'}
      \ar[r]\ar[d] &
    (\mathscr{I}/\mathscr{I}^n)
      \otimes_{\mathcal{O}_Y}\mathcal{O}_{Y'}
      \ar[r]\ar[d] & 0 \\
    0 \ar[r] &
    \mathscr{I}'^n/\mathscr{I}'^{n+1}
      \ar[r] &
    \mathscr{I}'/\mathscr{I}'^{n+1}
      \ar[r] &
    \mathscr{I}'/\mathscr{I}'^n
      \ar[r] & 0
  }
\]
en los cuales las filas son exactas (la primera por planitud
(0\textsubscript{I}, 6.1.2)) y las dos últimas flechas verticales son
biyectivas en virtud de (16.2.3, (ii)); de donde se sigue la conclusión.

\begin{remarks}[16.2.5]
\label{IV.16.2.5}
\begin{enumerate}
\item[(i)] El razonamiento de (16.2.2, (i)) sigue siendo aplicable cuando en
(16.2.1.1) se trata de morfismos de espacios localmente anillados
($\mathbf{Err}_{\mathrm{II}}$, (1.8.2)).
\item[(ii)] En (16.2.2, (ii)), la conclusión ya no es necesariamente válida
cuando se supone solamente que $v$ y $f$ son morfismos de preesquemas
($f$ satisface la condición de (16.1.1)). Por ejemplo (con las notaciones de la
demostración de (16.2.2, (ii))), puede ocurrir que $\mathfrak{I}=0$ pero que
el núcleo $\mathfrak{I}'$ de $A'\to B'=B\otimes_A A'$ no sea nulo y que
$B'\neq 0$, en cuyo caso se tiene $Y^{(n)}=Y$ para todo $n$, pero
$Y'^{(n)}\neq Y'$. Se obtiene un ejemplo de este hecho tomando
$A=\mathbb{Z}$, $B=\mathbb{Q}$,
$A'=\prod_{h=1}^{\infty}(\mathbb{Z}/m^h\mathbb{Z})$ donde $m>1$.
\end{enumerate}
\end{remarks}

\begin{env}[16.2.6]
\label{IV.16.2.6}
Consideremos el caso particular del diagrama (16.2.1.1) en que $X'=X$, $v$ es
la identidad, $X$ es un preesquema, $Y$ un subpreesquema de $X$, $Y'$ un
subpreesquema de $Y$, y $f$, $u$ y $f'=f\circ u$ las inyecciones canónicas; el
dihomomorfismo (16.2.1.3) da, por tensorización con
$\mathcal{O}_{Y'}$ sobre $\rho^*(\mathcal{O}_Y)$, un homomorfismo de
$\mathcal{O}_{Y'}$-álgebras graduadas
\[
  u^*(\mathscr{G}r_\bullet(f))
  \longrightarrow\mathscr{G}r_\bullet(f')
  \tag{16.2.6.1}
  \label{IV.16.2.6.1}
\]
Por otra parte, se identifica $\mathcal{O}_Y$ con
$\psi^*(\mathcal{O}_X)/\mathscr{I}_f$, y $\mathcal{O}_{Y'}$ con
$\rho^*(\mathcal{O}_Y)/\mathscr{I}_u$; como $\rho^*$ es un funtor exacto,
se tiene
$\rho^*(\mathcal{O}_Y)=
\rho^*(\psi^*(\mathcal{O}_X))/\rho^*(\mathscr{I}_f)=
\psi'^*(\mathcal{O}_X)/\rho^*(\mathscr{I}_f)$, y como
$\mathcal{O}_{Y'}$ se identifica además con
$\psi'^*(\mathcal{O}_X)/\mathscr{I}_{f'}$, se ve que se tiene
$\mathscr{I}_u=\mathscr{I}_{f'}/\rho^*(\mathscr{I}_f)$. Se deduce para
todo entero $n$ un homomorfismo canónico
$\mathscr{I}_{f'}^n/\mathscr{I}_{f'}^{n+1}\to
\mathscr{I}_u^n/\mathscr{I}_u^{n+1}$, de donde se obtiene un homomorfismo canónico de
$\mathcal{O}_{Y'}$-álgebras graduadas
\[
  \mathscr{G}r_\bullet(f')
  \longrightarrow\mathscr{G}r_\bullet(u).
  \tag{16.2.6.2}
  \label{IV.16.2.6.2}
\]
\end{env}

\begin{proposition}[16.2.7]
\label{IV.16.2.7}
Sean $X$ un preesquema, $Y$ un subpreesquema de $X$, $Y'$ un
subpreesquema de $Y$, $j:Y'\to Y$ la inyección canónica. Se tiene entonces una sucesión
exacta de haces conormales ($\mathcal{O}_{Y'}$-módulos)
\[
  \xymatrix{
    j^*(\mathscr{N}_{Y/X}) \ar[r] &
    \mathscr{N}_{Y'/X} \ar[r] &
    \mathscr{N}_{Y'/Y} \ar[r] & 0
  }
  \tag{16.2.7.1}
  \label{IV.16.2.7.1}
\]
donde las flechas son las componentes de grado 1 de los homomorfismos canónicos
(16.2.6.1) y (16.2.6.2).
\end{proposition}

Como la cuestión es local, basta considerar el caso en que
$X=\operatorname{Spec}(A)$, $Y=\operatorname{Spec}(A/\mathfrak{I})$,
$Y'=\operatorname{Spec}(A/\mathfrak{K})$, $\mathfrak{I}$ y $\mathfrak{K}$
son dos ideales de $A$ tales que $\mathfrak{I}\subset\mathfrak{K}$; todo
se reduce entonces a ver
\oldpage[IV-4]{14}
que la sucesión de homomorfismos canónicos
$\mathfrak{I}/\mathfrak{K}\mathfrak{I}\to
\mathfrak{K}/\mathfrak{K}^2\to
(\mathfrak{K}/\mathfrak{I})/(\mathfrak{K}/\mathfrak{I})^2\to 0$
es exacta, lo que es inmediato, puesto que la imagen de
$\mathfrak{I}/\mathfrak{K}\mathfrak{I}$ en
$\mathfrak{K}/\mathfrak{K}^2$ es
$(\mathfrak{I}+\mathfrak{K}^2)/\mathfrak{K}^2$ y que
$(\mathfrak{K}/\mathfrak{I})/(\mathfrak{K}/\mathfrak{I})^2$ se identifica con
$\mathfrak{K}/(\mathfrak{I}+\mathfrak{K}^2)$.

Es fácil dar ejemplos en los que la sucesión (16.2.7.1), prolongada a la izquierda
por un 0, deja de ser exacta; con las notaciones precedentes, basta
tomar $A=k[T]$, $\mathfrak{I}=AT^2$, $\mathfrak{K}=AT$, pues se tiene entonces
$(\mathfrak{I}+\mathfrak{K}^2)/\mathfrak{K}^2=0$ y
$\mathfrak{I}/\mathfrak{K}\mathfrak{I}\neq 0$. Véanse, sin embargo, (16.9.13) y
(19.1.5) para casos útiles en los que la sucesión prolongada sigue siendo exacta.
